% ============================================================
%  کتاب‌شناسی و منابع
% ============================================================
\section*{کتاب‌شناسی و منابع}
\addcontentsline{toc}{section}{کتاب‌شناسی و منابع}
\label{sec:bibliography}

\begin{tcolorbox}[wavebox, title={ساختار کتاب‌شناسی}]
منابع در پنج دسته ارائه شده‌اند: (۱) آثار کلاسیک بنیادین، (۲) آثار تحلیلی معاصر، (۳) سنت قاره‌ای، (۴) سنت اسلامی–ایرانی، (۵) مطالعات تطبیقی و منابع فارسی.
\end{tcolorbox}

\vspace{0.5cm}

% ─── بخش اول: آثار کلاسیک ──────────────────────────────────
\subsection*{الف. آثار کلاسیک بنیادین}

\begin{enumerate}[leftmargin=2em, label=\arabic*.]

\item \lr{Aristotle. \textit{Politics}. Trans. C.D.C. Reeve. Indianapolis: Hackett, 1998.}
\item \lr{Augustine of Hippo. \textit{De Libero Arbitrio} [\textit{On Free Will}]. In P. King (trans.), \textit{Against the Academicians and The Teacher}. Indianapolis: Hackett, 1995.}
\item \lr{Hobbes, T. (1651/1994). \textit{Leviathan}. Ed. E. Curley. Indianapolis: Hackett.}
\item \lr{Locke, J. (1689/1988). \textit{Two Treatises of Government}. Ed. P. Laslett. Cambridge: Cambridge University Press.}
\item \lr{Rousseau, J.-J. (1762/1968). \textit{The Social Contract}. Trans. M. Cranston. Harmondsworth: Penguin.}
\item \lr{Kant, I. (1785/1998). \textit{Groundwork for the Metaphysics of Morals}. Trans. M. Gregor. Cambridge: Cambridge University Press.}
\item \lr{Mill, J.S. (1859/1978). \textit{On Liberty}. Ed. E. Rapaport. Indianapolis: Hackett.}
\item \lr{Hegel, G.W.F. (1820/1991). \textit{Elements of the Philosophy of Right}. Trans. H.B. Nisbet. Cambridge: Cambridge University Press.}
\item \lr{Marx, K. (1844/2007). \textit{Economic and Philosophic Manuscripts of 1844}. Trans. M. Milligan. Mineola: Dover.}
\item \lr{Constant, B. (1819). "The Liberty of Ancients Compared with That of Moderns." In B. Fontana (Ed.), \textit{Political Writings}. Cambridge: Cambridge University Press, 1988.}

\end{enumerate}

% ─── بخش دوم: آثار تحلیلی معاصر ────────────────────────────
\subsection*{ب. آثار تحلیلی معاصر (قرن ۲۰–۲۱)}

\begin{enumerate}[leftmargin=2em, label=\arabic*., resume]

\item \lr{Berlin, I. (1969). \textit{Four Essays on Liberty}. Oxford: Oxford University Press.} [\rl{«دو مفهوم از آزادی» مقاله‌ی کلیدی این مجموعه است}]
\item \lr{Rawls, J. (1971). \textit{A Theory of Justice}. Cambridge, MA: Harvard University Press.}
\item \lr{Rawls, J. (1993). \textit{Political Liberalism}. New York: Columbia University Press.}
\item \lr{Nozick, R. (1974). \textit{Anarchy, State, and Utopia}. New York: Basic Books.}
\item \lr{Hayek, F.A. (1960). \textit{The Constitution of Liberty}. Chicago: University of Chicago Press.}
\item \lr{MacCallum, G.C. (1967). "Negative and Positive Freedom." \textit{Philosophical Review}, 76(3), 312--334.}
\item \lr{Taylor, C. (1979). "What's Wrong with Negative Liberty?" In A. Ryan (Ed.), \textit{The Idea of Freedom}. Oxford: Oxford University Press.}
\item \lr{Pettit, P. (1997). \textit{Republicanism: A Theory of Freedom and Government}. Oxford: Clarendon Press.}
\item \lr{Pettit, P. (2014). \textit{Just Freedom: A Moral Compass for a Complex World}. New York: W.W. Norton.}
\item \lr{Skinner, Q. (1998). \textit{Liberty before Liberalism}. Cambridge: Cambridge University Press.}
\item \lr{Sen, A. (1999). \textit{Development as Freedom}. New York: Anchor Books.} [\rl{\textit{توسعه به مثابه آزادی}. ترجمه: محمدسعید نوری نائینی. تهران: نشر کویر}]
\item \lr{Nussbaum, M. (2011). \textit{Creating Capabilities: The Human Development Approach}. Cambridge, MA: Harvard University Press.}
\item \lr{Sandel, M. (1982). \textit{Liberalism and the Limits of Justice}. Cambridge: Cambridge University Press.}
\item \lr{Dworkin, R. (2000). \textit{Sovereign Virtue: The Theory and Practice of Equality}. Cambridge, MA: Harvard University Press.}
\item \lr{Raz, J. (1986). \textit{The Morality of Freedom}. Oxford: Clarendon Press.}
\item \lr{Fricker, M. (2007). \textit{Epistemic Injustice: Power and the Ethics of Knowing}. Oxford: Oxford University Press.}

\end{enumerate}

% ─── بخش سوم: سنت قاره‌ای ──────────────────────────────────
\subsection*{ج. سنت قاره‌ای و پسامدرن}

\begin{enumerate}[leftmargin=2em, label=\arabic*., resume]

\item \lr{Sartre, J.-P. (1946/2007). \textit{Existentialism is a Humanism}. Trans. C. Macomber. New Haven: Yale University Press.}
\item \lr{Arendt, H. (1958). \textit{The Human Condition}. Chicago: University of Chicago Press.}
\item \lr{Arendt, H. (1961). \textit{Between Past and Future}. New York: Viking.}
\item \lr{Heidegger, M. (1927/1962). \textit{Being and Time}. Trans. J. Macquarrie \& E. Robinson. New York: Harper \& Row.}
\item \lr{Foucault, M. (1970). \textit{The Order of Discourse} [Inaugural lecture]. Trans. R. Young. London: Routledge, 1981.}
\item \lr{Foucault, M. (1977). \textit{Discipline and Punish}. Trans. A. Sheridan. New York: Pantheon.}
\item \lr{Butler, J. (1990). \textit{Gender Trouble: Feminism and the Subversion of Identity}. New York: Routledge.}
\item \lr{Habermas, J. (1984). \textit{The Theory of Communicative Action} (Vol. 1). Trans. T. McCarthy. Boston: Beacon Press.}
\item \lr{Taylor, C. (1989). \textit{Sources of the Self: The Making of the Modern Identity}. Cambridge, MA: Harvard University Press.}
\item \lr{Honneth, A. (1992/1995). \textit{The Struggle for Recognition}. Trans. J. Anderson. Cambridge: Polity Press.}
\item \lr{Pateman, C. (1988). \textit{The Sexual Contract}. Stanford: Stanford University Press.}
\item \lr{Fanon, F. (1961/2004). \textit{The Wretched of the Earth}. Trans. R. Philcox. New York: Grove Press.}
\item \lr{Spivak, G.C. (1988). "Can the Subaltern Speak?" In C. Nelson \& L. Grossberg (Eds.), \textit{Marxism and the Interpretation of Culture}. Urbana: University of Illinois Press.}

\end{enumerate}

% ─── بخش چهارم: سنت اسلامی–ایرانی ─────────────────────────
\subsection*{د. سنت اسلامی–ایرانی}

\begin{enumerate}[leftmargin=2em, label=\arabic*., resume]

\item ابن خلدون (۱۳۸۰). \textit{مقدمه‌ی ابن خلدون}. ترجمه: محمدپروین گنابادی. تهران: علمی و فرهنگی.

\item نظام‌الملک (۱۳۸۵). \textit{سیاست‌نامه}. ویرایش: محمد قزوینی. تهران: زوار.

\item غزالی (۱۳۸۱). \textit{نصیحة الملوک}. ویرایش: جلال‌الدین همایی. تهران: انجمن آثار ملی.

\item طباطبایی، سیدجواد (۱۳۸۵). \textit{ابن‌خلدون و علوم اجتماعی}. تهران: ققنوس.

\item طباطبایی، سیدجواد (۱۳۷۵). \textit{زوال اندیشه‌ی سیاسی در ایران}. تهران: کویر.

\item شایگان، داریوش (۱۳۸۱). \textit{افسون‌زدگی جدید}. تهران: فرزان.

\item سروش، عبدالکریم (۱۳۸۰). \textit{بسط تجربه‌ی نبوی}. تهران: صراط.

\item شریعتی، علی (۱۳۷۵). \textit{اسلام‌شناسی}. تهران: الهام.

\item \lr{An-Na'im, A. (1990). \textit{Toward an Islamic Reformation: Civil Liberties, Human Rights, and International Law}. Syracuse: Syracuse University Press.}
\item \lr{Soroush, A. (2000). \textit{Reason, Freedom, and Democracy in Islam}. Trans. \& ed. M. Sadri \& A. Sadri. Oxford: Oxford University Press.}

\end{enumerate}

% ─── بخش پنجم: فارسی و تطبیقی ──────────────────────────────
\subsection*{ه. منابع فارسی و مطالعات تطبیقی}

\begin{enumerate}[leftmargin=2em, label=\arabic*., resume]

\item آقایی، بهمن (۱۳۹۰). \textit{فرهنگ علوم سیاسی}. تهران: نشر علم.

\item بشیریه، حسین (۱۳۸۴). \textit{لیبرالیسم و محافظه‌کاری}. تهران: نشر نی.

\item بشیریه، حسین (۱۳۸۰). \textit{آموزش دانش سیاسی}. تهران: نشر نگاه معاصر.

\item زیبا کلام، صادق (۱۳۷۷). \textit{سنت و مدرنیته}. تهران: روزنه.

\item کاتوزیان، محمدعلی (۱۳۹۴). \textit{ایرانیان: دوران باستان تا دوره‌ی معاصر}. ترجمه: عبدالله کوثری. تهران: نشر مرکز.

\item \lr{Carter, I., Kramer, M., \& Steiner, H. (Eds.). (2007). \textit{Freedom: A Philosophical Anthology}. Oxford: Blackwell.} [\rl{مجموعه‌ای جامع از متون اصلی}]
\item \lr{Miller, D. (Ed.). (1991). \textit{Liberty}. Oxford: Oxford University Press.}

\end{enumerate}

\vspace{1cm}
\begin{center}
\begin{tcolorbox}[enhanced, width=10cm, colback=BgWarm,
  colframe=AccentGold, rounded corners, drop shadow]
\centering
{\large\bfseries\color{PrimaryDeep}پایان مقاله}\\[0.3cm]
{\small\color{NeutralDark}
نوشته‌ی \textbf{مهدی سالم}\\
۳ اسفند ۱۴۰۴ — ۲۲ فوریه ۲۰۲۶\\[0.2cm]
\lr{\url{MahdiSalem.com}} % FIXED: \lr{\url{...}}
}
\end{tcolorbox}
\end{center}
